% paralelos.tex

\chapter{Paralelos y polos}

En este apartado veremos dos elementos que se complementan, por lo que aparecen juntos,
aunque no es necesario. Se trata de los paralelos de una esfera y sus los polos norte
y sur (los puntos con $\theta=\qty{0}{\degree}$ y $\theta=\qty{180}{\degree}$,
respectivamente.

Como en todos los casos, para representar uno o más polos y/o paralelos en la esfera, se
utiliza el comando principal definido para Lua\LaTeX :
\begin{center}
  {\small\texttt{\textbackslash dibujaTikZEsfera[true]}}
  
  {\small\texttt{\textbackslash dibujaTikZEsfera[false]}}
\end{center}

Recordamos que, con \texttt{true} se indica que la esfera es total o parcialmente
transparente (translúcida) y con \texttt{false}, que es opaca.
De la misma manera que con los meridianos, si no se indica nada, se supone que es opaca:
\begin{center}
  {\small\texttt{\textbackslash dibujaTikZEsfera}}
\end{center}

Con el comando anterior se leen, además de \texttt{esfera}, oros campos o subtablas,
como \texttt{'observador'}, \texttt{'paralprops'} y \texttt{'paralelos} para representar
paralelos. Todo esto, excepto \texttt{'paralelos'} es igual que con los meridianos,
aunque lo repetiremos aquí:

\begin{enumerate}
\item El campo \texttt{'observador'} define la vista de un observador situado muy lejos
  de la esfera, definida por las coordenadas angulares esféricas polar $\theta$ y
  azimutal $\phi$ en grados:
  \begin{itemize}
    \setlist[itemize]{noitemsep}\vspace*{-1ex}
  \item \texttt{thetaD}: Ángulo polar de observación de la esfera, en grados.
  \item \texttt{phiD}: Ángulo azimutal de observación, en grados.
  \end{itemize}
\item El campo \texttt{'paralprops'} define propiedades visuales por defecto del dibujo
  de los meridianos, y otras de la parte \emph{visible} y la \emph{invisible} de los
  meridianos (si la esfera fuera opaca):
  \begin{itemize}
    \setlist[itemize]{noitemsep}\vspace*{-1ex}
  \item \texttt{loops}: Cuántas veces se repite el bucle de trazado de cada meridano
    (de cuantos puntos consta cada meridiano).
  \item \texttt{color\_vis}: Color reservado para el dibujo de la parte \emph{visible}.
  \item \texttt{lw\_vis}: Grosor de línea de la parte \emph{visible}.
  \item \texttt{color\_novis}: Color reservado para el dibujo de la parte
    \emph{invisible}.
  \item \texttt{lw\_novis}: Grosor de línea de la parte \emph{invisible} (dimensión
    \LaTeX , preferiblemente en puntos (\unit{pt}))
  \end{itemize}  
\item El campo \texttt{'paralelos'}, define el ángulo polar $\theta$ que define el
  paralelo en grados:
  \begin{itemize}
    \setlist[itemize]{noitemsep}\vspace*{-1ex}
  \item \texttt{thetaD}: Ángulo azimutal que define el meridiano, en grados (número)
  \end{itemize}
  Además del valor anterior, se puede destacar más o menos uno o más paralelos
  con respecto a los demás, añadiendo todos o alguno de los valores del
  campo \texttt{'paralprops'}, excepto \texttt{'loops'}.
\item El campo \texttt{'polprops'} define las características de los puntos que
  representan los polos norte y sur de la esfera:
  \begin{itemize}
  \item \texttt{'color\_vis'}: Color de los polos visibles (cadena de texto).
  \item \texttt{'radio\_vis'}: Radio en dimensiones \LaTeX , preferiblemente \unit{pt}
    de los polos visibles (cadena de texto).
  \item \texttt{'color\_novis'}: Color de los polos invisibles (cadena de texto).
  \item \texttt{'radio\_novis'}: Radio en dimensiones \LaTeX , preferiblemente \unit{pt}
    de los polos invisibles (cadena de texto). 
  \end{itemize}
\end{enumerate}

Veamos uno por uno los campos en los ficheros que mostraremos en los siguientes
ejemplos, \texttt{esf\_paral\_01.lua} y \texttt{esf\_paral\_02.lua}, obviando
los campos
\begin{center}\texttt{'esfera'}, \texttt{smbresfera} y \texttt{'observador'}\end{center}

\vspace{1ex}
CAMPO \texttt{'polprops'}:

\begin{tabular}{|@{\hspace{-0.8em}}c|}\hline
  \vphantom{\rule{0.4pt}{3ex}}
  \begin{lstlisting}[style=myLuastyle]
    polprops = {
      color_vis = "black!70", radio_vis="1pt",
      color_novis = "black!40", radio_novis="1pt",
    },
  \end{lstlisting}
  \\[1ex]
  \hline
\end{tabular}


\vspace{1ex}
CAMPO \texttt{'paralprops'}:

\begin{tabular}{|@{\hspace{-0.8em}}c|}\hline
  \vphantom{\rule{0.4pt}{3ex}}
  \begin{lstlisting}[style=myLuastyle]
    paralprops = {
      loops = 240, color_vis = "black!60", lw_vis="0.4pt",
      color_novis = "black!30", lw_novis="0.3pt",
    },
  \end{lstlisting}
  \\[1ex]
  \hline
\end{tabular}

\vspace{1ex}
CAMPO \texttt{'paralelos'}:

\begin{tabular}{|@{\hspace{-0.8em}}c|}\hline
  \vphantom{\rule{0.4pt}{3ex}}
  \begin{lstlisting}[style=myLuastyle]
   meridianos = {
     {phiD = 30}, {phiD = 60},
     {phiD = 90}, {phiD = 120}, {phiD = 150},
     {phiD = 180}, {phiD = -30}, {phiD = -60},
     {phiD = -90}, {phiD = -120},{phiD = -150}
   },
  \end{lstlisting}
  \\[1ex]
  \hline
\end{tabular}

En el fichero \texttt{esf\_paral\_01.lua} se define una esfera plana:

\vspace{1ex}
\begin{tabular}{|@{\hspace{-1em}}c|}\hline
  \vphantom{\rule{0.4pt}{3ex}}
  \begin{lstlisting}[basicstyle=\ttfamily\footnotesize, language=TeX]
    esfera = {radio=2.0, draw="black!25", lw="0.6pt", fill="black!8", opac=1.0,},
  \end{lstlisting}
  \\[1ex]
  \hline
\end{tabular}

Mientras que en \texttt{esf\_paral\_02.lua} está sombreada

\vspace{1ex}
\begin{tabular}{|@{\hspace{-1em}}c|}\hline
  \vphantom{\rule{0.4pt}{3ex}}
  \begin{lstlisting}[basicstyle=\ttfamily\footnotesize, language=TeX]
   esfera = {radio=2.0, draw="black!25", lw="0.6pt", fill="black!8", opac=1.0,},
   smbresfera = {ballcolor="white", opac=0.4},
  \end{lstlisting}
  \\[1ex]
  \hline
\end{tabular}

\begin{figure}[ht]
  \centering
  \begin{minipage}{0.45\textwidth}
    \cargaDatos{datoslua/esf_paral_01.lua}
    \begin{tikzpicture}[scale=1.0,]
      \dibujaTikzEsfera
    \end{tikzpicture}
  \end{minipage}
  \hfill
  \begin{minipage}{0.45\linewidth}
    \cargaDatos{datoslua/esf_paral_02.lua}
    \centering
    \begin{tikzpicture}[scale=1.0,]
      \dibujaTikzEsfera
    \end{tikzpicture}
  \end{minipage}
  \caption{Esferas opacas, plana y sombreada.}
\end{figure}

\begin{figure}[ht]
  \centering
  \begin{minipage}{0.45\textwidth}
    \cargaDatos{datoslua/esf_paral_01.lua}
    \begin{tikzpicture}[scale=1.0,]
      \dibujaTikzEsfera[true]
    \end{tikzpicture}
  \end{minipage}
  \hfill
  \begin{minipage}{0.45\linewidth}
    \cargaDatos{datoslua/esf_paral_02.lua}
    \centering
    \begin{tikzpicture}[scale=1.0,]
      \dibujaTikzEsfera[true]
    \end{tikzpicture}
  \end{minipage}
  \caption{Esferas transparentes, plana y sombreada.}
\end{figure}


%\vspace{1ex}
%En el fichero \texttt{esf\_merpar\_01.lua} representamos meridianos y paralelos
%en la esfera:
%
%\begin{figure}[ht]
%  \centering
%  \begin{minipage}{0.45\textwidth}
%    \cargaDatos{datoslua/esf_merpar_01.lua}
%    \begin{tikzpicture}[scale=1.0,]
%      \dibujaTikzEsfera
%    \end{tikzpicture}
%  \end{minipage}
%  \hfill
%  \begin{minipage}{0.45\linewidth}
%    \cargaDatos{datoslua/esf_merpar_01.lua}
%    \centering
%    \begin{tikzpicture}[scale=1.0,]
%      \dibujaTikzEsfera[true]
%    \end{tikzpicture}
%  \end{minipage}
%  \caption{Esferas opaca y translúcida con varios meridianos y paralelos.}
%\end{figure}


%\section{Segmentos de paralelos}
%
%\begin{figure}[ht]
%  \centering
%  \begin{minipage}{0.45\textwidth}
%    \cargaDatos{datoslua/esf_sgmpar_01.lua}
%    \begin{tikzpicture}[scale=1.0,]
%      \dibujaTikzEsfera
%    \end{tikzpicture}
%  \end{minipage}
%  \hfill
%  \begin{minipage}{0.45\linewidth}
%    \cargaDatos{datoslua/esf_sgmpar_02.lua}
%    \centering
%    \begin{tikzpicture}[scale=1.0,]
%      \dibujaTikzEsfera[true]
%    \end{tikzpicture}
%  \end{minipage}
%    \caption{Esferas opacas con un segmento paralelo. La de la izquierda es plana.
%      La de la derecha está sombreada.}
%\end{figure}













%%% Local Variables:
%%% coding: utf-8
%%% mode: latex
%%% TeX-engine: luatex
%%% TeX-master: "../elementos_esfera.tex"
%%% End:


